\section{Conclusion}
\label{sec:conclusion}

ViPragSent addresses two questions: how a frozen label-wise system compares with complete best-tuned baselines, and how aggregate performance should be read alongside label-specific leadership. On the fixed six-label evaluation, ViPragSent-Final reaches 84.3883 macro pragmatic F1, exceeding Vistral-7B by 1.5633 points. The result is an aggregate lead, not universal dominance: irony, idiom/figurative language, and code-switching remain narrowly below their respective leaders.

The contribution is a provenance-bounded evaluation resource, a documented development-only selection and freeze protocol, and traceable comparison artifacts rather than a claim of a universally superior architecture. The highest-priority next steps are controlled component ablations under the frozen protocol, independent reproduction when checkpoints are available, and resolution of source authorization and safe-access arrangements. Stored predictions, configuration, and hashes support inspection of the recorded evidence, but not an independent full training reproduction.
